\documentclass[11pt]{article}
\usepackage[a4paper,margin=25mm]{geometry}
\usepackage{amsmath,amssymb,amsthm}
\usepackage[hidelinks]{hyperref}
\emergencystretch=2em
\newtheorem{theorem}{Theorem}
\title{One Newton Polygon, Two Frobenius Polynomials:\\A Disproof of Conjecture 00000005556}
\author{GitHub: gaochengzhecpu}
\date{}
\begin{document}
\maketitle
\begin{abstract}
Over $\mathbb F_3$, take the rank-one unit-root $F$-crystals with
Frobenius $+1$ and $-1$ on the free $\mathbb Z_3$-line. Their
rationalizations are actual linear automorphisms of $\mathbb Q_3$.
Both Newton polygons are the horizontal segment from $(0,0)$ to
$(1,0)$, and both crystals are pure of weight zero. Their Frobenius
characteristic polynomials are nevertheless $X-1$ and $X+1$.
Thus the Newton polygon does not determine the characteristic
polynomial, even at fixed rank and pure weight in the unit-root case.
Lean verifies the integral and rational automorphisms, scalar
extension compatibility, valuations, actual matrix characteristic
polynomials, geometric Newton segments and the universal
non-determination statement.
\end{abstract}

\paragraph{Reading of the source.}
The first assertion in both source languages is that the
characteristic polynomial is given by the Newton polygon. In
particular, this requires equal Newton polygons to give equal
characteristic polynomials. We refute this necessary uniqueness
condition, which already suffices to refute the stated conjecture.

The final phrase, saying that in the unit-root case something is
uniquely determined by pure weight, has an ambiguous subject. If
it refers to the characteristic polynomial, the same example also
refutes it. If it refers only to the Newton polygon, that weaker
claim is not refuted: all rank-one unit-root objects already have
the same horizontal Newton polygon. We do not claim otherwise.

\paragraph{The actual crystalline objects.}
Fix the base point $\operatorname{Spec}\mathbb F_3$. Here
$W(\mathbb F_3)=\mathbb Z_3$, its fraction field is
$\mathbb Q_3$, and the Witt Frobenius is the identity. Thus a free
rank-one $F$-crystal on this point is described by a free
$\mathbb Z_3$-module $M$ of rank one with an injective linear
Frobenius. It is unit-root when Frobenius is an automorphism of
this integral lattice. Its rationalization is an automorphism of
$M\otimes_{\mathbb Z_3}\mathbb Q_3$.

This is the specialization of the usual definition with a
Frobenius-semilinear automorphism: when the base Frobenius is the
identity, semilinearity is ordinary linearity. For the general
definition see Jacob Lurie, \emph{Lecture 26--Isocrystals},
Definition 1:
\href{https://www.math.ias.edu/~lurie/205notes/Lecture26-Isocrystals.pdf}{lecture notes}.
The standard identifications with Witt vectors and this
specialization are stated here as mathematical background.
Lean directly verifies all the resulting integral and rational
linear-algebra data over the actual $3$-adic rings and field.

Use the following two objects:
\[
 \begin{array}{c|c|c}
 &M&F\text{ on }M\text{ and its rationalization}\\\hline
 C_+&\mathbb Z_3&x\longmapsto x\\
 C_-&\mathbb Z_3&x\longmapsto -x.
 \end{array}
\]
Both maps and their inverses preserve the integral lattice. The
integral maps agree, under the inclusion
$\mathbb Z_3\hookrightarrow\mathbb Q_3$, with the displayed
rational maps. There is no quotient, reduction of coefficients, or
change of base in this comparison.

\begin{theorem}
The two unit-root $F$-crystals $C_+$ and $C_-$ have equal Newton
polygons and equal pure weight zero but distinct characteristic
polynomials. Their rational Frobenius operators are not conjugate
by a $\mathbb Q_3$-linear isomorphism.
\end{theorem}
\begin{proof}
The Frobenius matrices in the basis $1$ are respectively $[1]$
and $[-1]$, giving
\[
 \chi_+(X)=X-1,\qquad \chi_-(X)=X+1.
\]
They differ because $1\ne-1$ in the characteristic-zero field
$\mathbb Q_3$.

Let $v_3$ be the $3$-adic valuation with $v_3(3)=1$.
Both eigenvalues are units:
\[
 |1|_3=|-1|_3=1,\qquad v_3(1)=v_3(-1)=0.
\]
For a monic degree-one polynomial $c_0+c_1X$ with $c_0\ne0$,
its Newton polygon is the segment joining
$(0,v_3(c_0))$ and $(1,v_3(c_1))$. Thus the polygons of
$X-1$ and $X+1$ both equal
\[
 \{(s,0):0\le s\le1\}.
\]
The eigenvalues $1,-1$ are rational algebraic numbers, and every
complex embedding of their number field $\mathbb Q$ has absolute
value one on them. Since $3^{0/2}=1$, both objects are pure of
weight zero in the usual Weil-number sense. This is a stronger
counterexample than one obtained by changing rank or weight.

Finally, a linear isomorphism $f$ intertwining the two operators
would satisfy $-f(1)=f(1)$, hence $2f(1)=0$ and $f(1)=0$.
Injectivity would imply $1=0$, impossible.
\end{proof}

\paragraph{Why this refutes determination.}
A Newton polygon records coefficient valuations, or equivalently
the corresponding valuation slopes, rather than the full unit
parts of Frobenius eigenvalues. The above example preserves every
Newton datum while changing the unit $+1$ to $-1$. Therefore no
function of the Newton polygon alone can return the characteristic
polynomial for all crystals in the stated class. Even supplying
the pure weight zero in addition does not distinguish this pair.
This argument stays over the fixed finite base field
$\mathbb F_3$; it makes no assertion about classification after
extension to an algebraically closed base.

\paragraph{Formal verification and scope.}
The Lean structure \texttt{RankOneCrystal} contains an actual
$\mathbb Z_3$-linear equivalence of the integral line, an actual
$\mathbb Q_3$-linear equivalence of the rational line, and a proof
of scalar extension compatibility. The two examples instantiate
all three fields. The theorem
\begin{center}\texttt{matrix\_represents\_frobenius}\end{center}
proves that the actual one-by-one Mathlib matrix acts as the
rational Frobenius on every vector. Characteristic polynomials
are Mathlib's \texttt{Matrix.charpoly}. The valuations and norms
are those of Mathlib's actual $3$-adic field, and the Newton
polygon is a standard real geometric \texttt{segment} joining
the two coefficient-valuation points. The valuation calculation
does not substitute integers or a hard-coded slope table for
$3$-adic values.

\texttt{PureWeightZero} requires a rational Frobenius eigenvalue
and modulus one under every complex embedding of $\mathbb Q$;
this is exactly the usual purity condition for the rational
eigenvalues in our counterexample. It is not intended as a
definition of all possibly irrational pure crystals.
The principal negations are
\begin{center}
\texttt{not\_newton\_determines\_characteristic\_polynomial},\\
\texttt{not\_pure\_weight\_determines\_characteristic\_polynomial}.
\end{center}
These are universally quantified propositions over the actual
specialized crystalline objects; the first additionally requires
equal Newton polygons and unit-root conditions on both objects.
The second treats the stronger characteristic-polynomial reading
of the ambiguous final clause. Non-conjugacy is separately proved
by \texttt{no\_frobenius\_conjugacy}.

The Witt-vector identification and equivalence of this specialized
linear-algebra description with $F$-crystals on the point are
paper background. No formal construction of the full crystalline
site or of Mathlib's general Witt-vector isocrystal class is
claimed. The finite-field identity $x^3=x$ on $\mathbb F_3$ is
also checked directly in Lean.

\paragraph{Source and reproduction.}
The preserved bilingual source is \texttt{SOURCE.md}; the
provenance file records the repository commit and source hash.
The project pins Lean 4.19.0 and Mathlib revision
\begin{center}\texttt{c44e0c8ee63ca166450922a373c7409c5d26b00b}.\end{center}
From the \texttt{lean} directory run \texttt{lake build}, then
\begin{center}\texttt{lake env lean -DwarningAsError=true Main.lean}.\end{center}
Principal axiom reports contain only \texttt{propext},
\texttt{Classical.choice}, and \texttt{Quot.sound}.
Development and self-review were performed by a delegated AI
agent; the coordinating agent reviews the result separately.
No external independent review is claimed.
\end{document}
